%================================================
\section{Memory-Driven Model}
\begin{frame}{Memory and Non-Markovian Dynamics}
	\vspace{-0.4cm}

	\begin{columns}[T]

		\column{0.48\textwidth}

		\begin{block}{Markovian Dynamics}
			\[
				P(X_{t+\Delta t}\mid X_t)
			\]

			\begin{itemize}
				\item Future depends only on the current state.
				\item No memory.
				\item Instantaneous response.
			\end{itemize}
			\vspace{-0.6cm}

			\[
				m\dot v(t)=-\gamma v(t)+\eta(t)
			\]
			\vspace{-0.8cm}

			\[
				K(t)=\gamma\delta(t)
			\]

		\end{block}

		\column{0.48\textwidth}

		\begin{block}{Non-Markovian Dynamics}
			\[
				P(X_{t+\Delta t}\mid X_t,X_{t-\tau},\ldots)
			\]

			\begin{itemize}
				\item Future depends on past history.
				\item Memory effects.
				\item Persistent correlations.
			\end{itemize}
			\vspace{-0.8cm}

			\[
				m\dot v(t)
				=
				-\int_0^t K(t-s)v(s)\,ds
				+\eta(t)
			\]
			\vspace{-0.8cm}

			\[
				K(t)=K_0e^{-t/\tau_m}
			\]

		\end{block}

	\end{columns}

	\vspace{0.3cm}

	\begin{alertblock}{Motivation}
		Finite VACF and OCF decay indicate that information about past
		velocity and orientation persists for a finite time, motivating
		a non-Markovian description with memory kernels.
	\end{alertblock}

\end{frame}
% % %================================================
% % \begin{frame}{From Markovian to Non-Markovian Dynamics}
% % 	\vspace{-0.4cm}

% % 	\begin{columns}[T]

% % 		%---------------------------------------
% % 		\column{0.48\textwidth}

% % 		\begin{block}{Markovian Dynamics}

% % 			Future motion depends only on the current state.
% % 			% \vspace{-0.2cm}
% % 			\vspace{-0.2cm}

% % 			\[
% % 				P(X_{t+\Delta t}|X_t)
% % 			\]

% % 			% \vspace{-0.5cm}
% % 			% \[
% % 			% 	m\dot v(t)
% % 			% 	=
% % 			% 	-\gamma v(t)+\eta(t)
% % 			% \]

% % 			\begin{itemize}
% % 				\item No memory.
% % 				\item Instantaneous response.
% % 				\item Brownian diffusion.
% % 			\end{itemize}

% % 		\end{block}

% % 		%---------------------------------------
% % 		\column{0.48\textwidth}

% % 		\begin{alertblock}{Non-Markovian Dynamics}

% % 			Future motion depends on past history.
% % 			% \vspace{-0.2cm}
% % 			\vspace{0.2cm}

% % 			\[
% % 				P(X_{t+\Delta t}|X_t,X_{t-\tau},\ldots)
% % 			\]

% % 			% \vspace{-0.7cm}

% % 			% \[
% % 			% 	m\dot v(t)
% % 			% 	=
% % 			% 	-\int_0^t K(t-s)v(s)\,ds
% % 			% 	+\eta(t)
% % 			% \]

% % 			\begin{itemize}
% % 				\item Memory effects.
% % 				\item History-dependent motion.
% % 				\item Generalized Langevin Equation.
% % 			\end{itemize}

% % 		\end{alertblock}

% % 	\end{columns}

% % 	\vspace{0.15cm}

% % 	\begin{alertblock}{Motivation}

% % 		The observed self-avoiding trajectories suggest that past motion influences future motion, motivating a non-Markovian description.

% % 	\end{alertblock}

% % \end{frame}
% %================================================

% \begin{frame}{Memory Kernels and Correlations}
% 	\vspace{-0.3cm}

% 	\begin{columns}[T]

% 		%---------------------------------------
% 		\column{0.48\textwidth}

% 		\begin{block}{Markovian Limit}

% 			\[
% 				K(t)=\gamma\delta(t)
% 			\]
% 			\vspace{-0.5cm}

% 			\[
% 				m\dot v(t)
% 				=
% 				-\gamma v(t)+\eta(t)
% 			\]

% 			\begin{itemize}
% 				\item Memoryless dynamics.
% 				\item Instantaneous response.
% 				\item Future depends only on the current state.
% 				\item Correlations decay rapidly.
% 			\end{itemize}

% 		\end{block}

% 		%---------------------------------------
% 		\column{0.48\textwidth}

% 		\begin{alertblock}{Finite Memory}

% 			\[
% 				K(t)=K_0 e^{-t/\tau_m}
% 			\]
% 			\vspace{-0.8cm}

% 			\[
% 				m\dot v(t)
% 				=
% 				-\int_0^t K(t-s)v(s)\,ds
% 				+\eta(t)
% 			\]

% 			\begin{itemize}
% 				\item Past states influence the present.
% 				\item Finite memory timescale $\tau_m$.
% 				\item Persistent temporal correlations.
% 				\item Future depends on motion history.
% 			\end{itemize}

% 		\end{alertblock}

% 	\end{columns}

% 	\vspace{0.2cm}

% 	\begin{alertblock}{Key Idea}

% 		A finite memory kernel produces persistent temporal correlations.

% 	\end{alertblock}

% \end{frame}

%================================================
\begin{frame}{Chemical Memory Field}
	\[
		\frac{\partial c}{\partial t}
		=
		D\nabla^2 c
		+
		\alpha G_{R_t}(x-X_t)
	\]

	\begin{itemize}
		\item Droplet continuously deposits chemical material.
		\item Deposited material diffuses slowly.
		\item The field retains information about past positions.
	\end{itemize}
	\begin{alertblock}{Key Idea}
		The chemical field evolves more slowly than the droplet motion,
		allowing information about previously visited locations to persist.
	\end{alertblock}
\end{frame}

%================================================
\begin{frame}{Non-Markovian Droplet Model}

	\begin{columns}[T]

		%================================================
		\column{0.48\textwidth}

		\begin{block}{Chemical Field (PDE)}

			\[
				\frac{\partial c}{\partial t}
				=
				D\nabla^2 c
				+
				\alpha G_{R_t}(\mathbf{x}-\mathbf{X}_t)
			\]

			\begin{itemize}
				\item Diffusion of chemical field.
				\item Moving droplet source.
				\item Stores memory.
			\end{itemize}

		\end{block}

		%================================================
		\column{0.48\textwidth}

		\begin{alertblock}{Droplet Motion (SDE)}

			\[
				\eta_t \dot{\mathbf X}_t
				=
				-\mathbf F_{\mathrm{chem}}
				+
				\sqrt{\sigma}\,\xi_t
			\]

			\vspace{-0.1cm}

			\[
				\mathbf F_{\mathrm{chem}}
				\propto
				\int_{\Omega}
				G_{R_t}
				\nabla c\,d\mathbf x
			\]

			\begin{itemize}
				\item Drift from chemical gradients.
				\item $-\nabla c$ gives self-repulsion.
				\item Noise drives fluctuations.
			\end{itemize}

		\end{alertblock}
	\end{columns}

	\vspace{0.15cm}

	\begin{alertblock}{Feedback Loop}

		\[
			\text{Motion}
			\rightarrow
			\text{Chemical Field}
			\rightarrow
			\text{Motion}
		\]

	\end{alertblock}

	\vspace{0.2cm}

	\rightcite{droplet2025}

\end{frame}

%================================================
\begin{frame}{Model Parameters and Notation}
	\vspace{-0.2cm}

	\begin{columns}[T]

		%================================================
		\column{0.58\textwidth}

		\begin{block}{Model Equations}

			\textbf{Chemical Field}

			\[
				\frac{\partial c}{\partial t}
				=
				D\nabla^2 c
				+
				\alpha G_{R_t}(\mathbf{x}-\mathbf{X}_t)
			\]

			\vspace{0.2cm}

			\textbf{Droplet Motion}
			\vspace{-0.4cm}

			\[
				\eta_t \dot{\mathbf X}_t
				=
				-
				\frac{\beta R_t^3}{\delta^3}
				\int_{\Omega}
				G_{R_t}(\mathbf{x}-\mathbf X_t)
				\nabla c(\mathbf{x},t)\,d\mathbf x
				+
				\sqrt{\sigma}\,\xi_t
			\]

		\end{block}

		\begin{alertblock}{Key Idea}
			Past droplet positions are encoded in the chemical field and influence future motion through chemical-gradient sensing.
		\end{alertblock}

		%================================================
		\column{0.38\textwidth}

		\begin{block}{Parameters}

			\centering

			\small

			\begin{tabular}{ll}
				\hline
				\textbf{Symbol} & \textbf{Meaning}         \\
				\hline
				$D$             & Diffusion coefficient    \\
				$\alpha$        & Source strength          \\
				$R_t$           & Droplet radius           \\
				$\eta_t$        & Friction coefficient     \\
				$\sigma$        & Noise strength           \\
				$\delta$        & Interaction length scale \\
				$\Omega$        & Computational domain     \\
				\hline
			\end{tabular}

		\end{block}

		\vspace{0.2cm}

	\end{columns}
	\rightcite{droplet2025}
\end{frame}
